\documentclass[handout]{beamer}
\usepackage{graphicx}
\usepackage{booktabs}
\usepackage{array}
\usepackage{listings}
\usepackage{xcolor}

\definecolor{codegreen}{rgb}{0,0.6,0}
\definecolor{codegray}{rgb}{0.5,0.5,0.5}
\definecolor{codepurple}{rgb}{0.58,0,0.82}
\definecolor{backcolour}{rgb}{0.95,0.95,0.92}

\lstset{
    backgroundcolor=\color{backcolour},   
    commentstyle=\color{codegreen},
    keywordstyle=\color{magenta},
    numberstyle=\tiny\color{codegray},
    stringstyle=\color{codepurple},
    basicstyle=\ttfamily\scriptsize,
    breakatwhitespace=false,         
    breaklines=true,                 
    captionpos=b,                    
    keepspaces=true,                 
    numbers=left,                    
    numbersep=5pt,                  
    showspaces=false,                
    showstringspaces=false,
    showtabs=false,                  
    tabsize=2
}

\title{Re-processing Perturb-Seq datasets from FASTQ}
\subtitle{Unified Pipeline and Infrastructure Challenges}
\date{2026-03-26}
\author{Kirill Tsukanov \\ Senior Full Stack Developer \& Data Engineer \\ \texttt{ktsukanov@ebi.ac.uk}}
\institute{PC Fortnightly Meeting}

\setbeamertemplate{navigation symbols}{}
\setbeamertemplate{footline}{%
  \begin{beamercolorbox}[wd=\paperwidth,ht=2.5ex,dp=1.5ex,center]{}
    \hspace{15pt}\usebeamerfont{author in head/foot}Unified FASTQ Re-analysis Pipeline
    \hfill
    \usebeamerfont{date in head/foot}\insertframenumber{} / 9
    \hspace{15pt}
  \end{beamercolorbox}%
}

\begin{document}

\begin{frame}
    \titlepage
\end{frame}

\begin{frame}{Why re-process from FASTQ?}
    \begin{itemize}
        \item \textbf{Consistency:} Most scPerturb datasets are provided as processed H5AD files, but they use different pipelines (Cell Ranger, STARsolo, kallisto) and different reference versions.
        \item \textbf{Modality Alignment:} Ensuring that the layout and metadata of \texttt{h5ad} files are consistent across studies. This reduces the need for per-dataset adjustments during manual curation.
        \item \textbf{Quality Control:} Direct access to raw reads would allow for uniform filtering and normalization across the entire Perturbation Catalogue.
        \item \textbf{Expectations:} When we talked to people, most of them \emph{expect} that data will be fully harmonised in the catalogue.
    \end{itemize}
\end{frame}

\begin{frame}{Data Sourcing: The ENA vs. SRA Dilemma}
    \begin{itemize}
        \item \textbf{The Issue:} For direct ENA downloads of single-cell data, it only hosts the cDNA reads for historical reasons.
        \item \textbf{The Consequence:} Technical reads (Cell Barcodes and UMIs) are missing from ENA FASTQ files.
        \item \textbf{The Solution:} We must fetch data from the \textbf{SRA (Sequence Read Archive)} using the SRA Toolkit.
        \item \textbf{Workflow:}
            \begin{enumerate}
                \item Query ENA API for run accessions (SRR).
                \item Use \texttt{prefetch} to download SRA files.
                \item Use \texttt{fasterq-dump} with \texttt{--include-technical}.
            \end{enumerate}
    \end{itemize}
\end{frame}

\begin{frame}{Infrastructure and Storage Constraints}
    \begin{itemize}
        \item \textbf{Limited Quota:} Only 20 TB allocated on the cluster storage.
        \item \textbf{Space-Heavy Process:} 
            \begin{itemize}
                \item \texttt{fasterq-dump} produces \textbf{uncompressed} FASTQ files.
                \item Uncompressed original FASTQ files can easily exceed the quota during the download phase, even before we get to the actual analysis.
            \end{itemize}
        \item \textbf{Storage Conscious Approach:} 
            \begin{itemize}
                \item Immediate \texttt{gzip} after every given \texttt{fasterq-dump} process completes.
                \item Nextflow \texttt{cleanup} enabled to purge intermediate results.
                \item Threaded downloads and processing to minimize "idle" storage usage.
            \end{itemize}
    \end{itemize}
\end{frame}

\begin{frame}{Unified Pipeline Architecture}
    \begin{itemize}
        \item \textbf{Framework:} Nextflow: scales well, works out of the box on the EBI Slurm cluster.
        \item \textbf{Engine:} \texttt{kb-python} (Bustools + Kallisto wrapper) for mapping.
        \item \textbf{Environment:} Singularity container, as it is the only solution that works properly on our SLURM cluster.
    \end{itemize}
    \vspace{0.5em}
    \begin{center}
        \textbf{Pipeline Flow:}
        \begin{enumerate}
            \item \textbf{Index Building:} Simultaneous cDNA (Ensembl) and KITE (CRISPR feature) indices.
            \item \textbf{Counting:} Parallel \texttt{kb count} for both modalities across all SRR runs.
            \item \textbf{Merging:} Per-SRR alignment of cells (matching barcodes).
            \item \textbf{Concatenation:} Efficient on-disk concatenation (\texttt{anndata.experimental}) of all runs into a single \texttt{experiment\_final.h5ad}.
        \end{enumerate}
    \end{center}
\end{frame}

\begin{frame}{Handling Dual Modalities (cDNA + CRISPR)}
    \begin{itemize}
        \item The pipeline runs two parallel workflows per SRR:
        \begin{enumerate}
            \item \textbf{Standard Workflow:} Maps cDNA reads to the transcriptome.
            \item \textbf{KITE Workflow:} Maps Feature Barcode reads to a guide whitelist (\texttt{features.tsv}).
        \end{enumerate}
        \item \textbf{Merging Logic:}
            \begin{itemize}
                \item We use the cDNA cell barcodes as the "truth" set.
                \item Guide counts are aligned to these barcodes.
                \item Final H5AD stores gene counts in \texttt{.X} and guide counts in \texttt{.obsm['guides']}.
            \end{itemize}
        \item \textbf{Auto-detection:} The pipeline automatically identifies R1 (barcode+UMI) and R2 (read) based on sequence length, simplifying run configuration.
    \end{itemize}
\end{frame}

\begin{frame}[fragile]{Pipeline Execution Example}
\begin{lstlisting}[language=bash]
nextflow run main.nf \
  -profile slurm,singularity \
  --fastq_dir $HPS_PATH/perturb_seq_fastq/SAMN40972597 \
  --outdir results \
  --chemistry 10xv3 \
  --transcriptome_fa $REF/Homo_sapiens.GRCh38.dna.primary_assembly.fa.gz \
  --gtf $REF/Homo_sapiens.GRCh38.111.gtf.gz \
  --features_tsv ./features.tsv
\end{lstlisting}
    \vspace{0.5em}
    \begin{itemize}
        \item Automated SLURM resource allocation.
        \item Built-in retry mechanism for transient cluster errors.
        \item \texttt{h5repack} with GZIP compression used for final archival.
    \end{itemize}
\end{frame}

\begin{frame}{Case Study: Nadig 2025 (Jurkat)}
    \begin{itemize}
        \item \textbf{Data Volume:} 
            \begin{itemize}
                \item Total Compressed FASTQ: $\sim$2 TB.
                \item Scale: 2,688 FASTQ files across multiple SRR runs.
            \end{itemize}
        \item \textbf{Performance (Nadig 2025 Jurkat):} 
            \begin{itemize}
                \item \textbf{Download:} $\sim$3 hours.
                \item \textbf{Processing:} $\sim$5 hours.
                \item \textbf{Archival:} $\sim$1 hour (compression + upload).
            \end{itemize}
        \item \textbf{Future Scaling:} 
            \begin{itemize}
                \item These times will be further improved.
                \item We aim for a very fast, resilient pipeline that scales for the entire catalogue.
            \end{itemize}
    \end{itemize}
\end{frame}

\begin{frame}{Next Steps}
    \begin{itemize}
        \item \textbf{Validation:} Perform correlation analysis between our re-processed counts and the authors' original H5AD files for a basic sanity check.
        \item \textbf{Scaling:} Extend processing to other large datasets.
        \item \textbf{Refinement:} Automate the generation of \texttt{features.tsv} from supplementary tables for more studies.
        \item \textbf{Integration:} Feed the unified H5ADs into our downstream pipelines: curation, differential expression, GSEA.
    \end{itemize}
\end{frame}

\end{document}
